\documentclass[11pt,a4paper,twocolumn]{article}

% === Core Packages ===
\usepackage[margin=2cm]{geometry}
\usepackage[utf8]{inputenc}
\usepackage[T1]{fontenc}
\usepackage{times}
\usepackage{amsmath,amssymb}
\usepackage{graphicx}
\usepackage{booktabs}
\usepackage{float}
\usepackage{enumitem}
\usepackage{tikz}
\usepackage{pgfplots}
\pgfplotsset{compat=1.18}
\usetikzlibrary{arrows.meta,shapes.geometric,positioning,calc,decorations.markings,fit,backgrounds}

% === Citation (IEEE style) ===
\usepackage[numbers,sort&compress]{natbib}
\bibliographystyle{IEEEtranN}

% === Hyperref ===
\usepackage[hidelinks]{hyperref}
\usepackage{xurl}

% === Settings ===
\setlength{\parindent}{1em}
\setlength{\parskip}{0.3em}

% === Title ===
\title{\textbf{탈중앙화 평판 네트워크:\\자연스러운 사회 조직을 위한 프레임워크}}
\author{Pavel Kudrna\\Prague, Czechia\\Draft v1 --- March 2026}
\date{}

\begin{document}
\maketitle

% ============================================================
% ABSTRACT
% ============================================================
\begin{abstract}
정의감---잘못이 바로잡히고 공로가 보상받는다는 확신---은 사회를 결속시키는 가장 강력한 힘이다. 중앙집권적 통치 기관이 권리를 집행하는 비용이 권리 자체의 가치를 초과하는 탓에 이 감각을 제공하지 못할 때, 사회적 결속은 침식된다. 현행 제도 구조는 중앙 계획이 자원 배분에 실패하는 것과 체계적으로 동일한 방식으로 상당한 부류의 분쟁을 해결하지 못한다. 즉 정보 비대칭, 피드백 루프의 부재, 그리고 의사결정자가 자신의 결정이 초래하는 결과로부터 단절되는 방식이다. 본 논문은 평판 소셜 네트워크(Reputation Social Network, RSN)를 제시한다. RSN은 탈중앙화 식별자(DID)를 기반으로 구축된 탈중앙화되고, 부패시킬 수 없으며, 검열에 저항하는 통신 계층으로서, 가명 참여자들이 현실 세계의 행동에 관한 평판 정보를 게시하고, 검증하며, 소비할 수 있게 한다. 핵심 메커니즘은 세 가지 설계 원칙에 기반한다. (1)~평판 데이터를 읽는 것은 저렴하지만 쓰는 것은 검증 가능한 노력을 요구하는 비대칭적 비용 구조. (2)~일관된 해싱(consistent hashing)에 기반한 비결정론적 검증자 선택 프로토콜로, 반복 비용의 점증을 통해 급진적 주장에 가격을 매기는 방식. (3)~민간 검증자가 자신이 보증하는 주장의 정확성에 자신이 축적한 평판을 거는 시장 주도형 평판 권위 모델. 그 결과로 생성되는 아키텍처는 중앙 조정 없이 실력주의(meritocracy)를 창발적으로 만들어내며, 기존의 국가 관리 체계로부터 점진적이고 되돌릴 수 있는 이행 경로를 가능하게 한다.
\end{abstract}

% ============================================================
% 1. INTRODUCTION
% ============================================================
\section{서론}

\subsection{중앙집권적 신뢰의 문제}

현대의 통치는 근본적인 가정에 기반한다. 즉 중앙집권적 기관이 낯선 사람들 사이의 신뢰를 대규모로 중개할 수 있다는 것이다. 법원은 분쟁을 판결한다. 등기소는 소유권을 증명한다. 규제 기관은 표준을 집행한다. 이러한 기관들은 정보가 희소하고, 소통이 느리며, 중앙 기구로의 권한 위임이 유일하게 실현 가능한 조정 메커니즘이던 시대에 설계되었다. 그러나 중앙집권적 통치는 중앙 계획과 동일한 구조적 이유로 실패한다. 정보 비대칭, 피드백 루프의 부재, 그리고 의사결정자가 자신의 결정이 초래하는 결과로부터 단절되는 것이다.

이 가정은 예컨대 제도적 중개의 비용이 분쟁의 가치를 초과할 때 무너진다. 임차한 아파트에 법적으로 요구되는 전기 안전 증명서가 없다는 것을 발견한 세입자를 생각해 보라. 이는 생명에 즉각적인 위험을 야기하는 상태다. 계약을 해지할 세입자의 법적 권리는 명백하다. 그러나 법원 시스템을 통해 그 권리를 집행하는 비용은, 잠재적인 법정 배상금을 포함하더라도, 보증금과 손해 배상액의 금전적 가치를 초과한다~\cite{czcourtfees}. 합리적인 대응은 손실을 감수하고 떠나는 것이다.

이는 병리적인 예외 사례가 아니다. 이는 피해는 실재하지만 금전적 이해가 제도적 집행이 경제적으로 합리적이 되는 문턱 아래로 떨어지는 상당한 부류의 분쟁에서 기본적으로 겪는 경험이다. 그 결과는 악의적 행위자에게 구조적으로 유리한 시스템이다. 이 문턱 아래의 규모로 타인에게 피해를 주는 자들은 처벌받지 않고 행동할 수 있는데, 정의를 구하는 비용이 피해를 입은 당사자에게 지워지기 때문이다.

위의 사례는 저자의 법적 환경---체코의 대륙법 체계---에서 가져온 것이다. 다른 법 전통에서는---판례 기반의 보통법, 관습법, 혼합 체계---관할권의 문화적, 절차적 특수성에 따라 정의가 서로 다른 방식과 정도로 실패한다. 독자들은 자신의 맥락에서 이에 상응하는 사례를 알아볼 가능성이 크다. 구조적으로 보편적인 것은 그 패턴이다. 즉 그 가치가 경제적으로 합리적인 집행의 문턱 아래로 떨어지는 분쟁은 피해를 입은 당사자가 손실을 떠안는 것으로 끝난다. RSN은 완벽한 정의를 약속하지 않는다. 다만 제도적 집행이 비경제적일 때 악의적 행위자가 누리는 구조적 이점을 줄일 뿐이다.

\subsection{기존 해법이 부족한 이유}

\textbf{탈중앙화 신원(DID) 프레임워크}~\cite{did2022}는 자기주권적 신원 관리를 위한 암호학적 메커니즘을 제공하지만, 행동 평판이라는 더 넓은 문제를 다루지 않고 주로 신원 검증에 초점을 맞춘다.

\textbf{소울바운드 토큰(SBT)}~\cite{weyl2022}은 평판이 양도 불가능하다는 통찰을 포착하지만, 확장성 제약이 있는 블록체인 인프라에 의존하며 정보 입력을 지배하는 경제적 유인을 다루지 않는다. 공로 토큰(예: Liberland)과 같은 관련 개념도 유사한 철학을 공유하지만 특정 관할권 프레임워크에 얽매여 있다.

\textbf{탈중앙화 자율 조직(DAO)}~\cite{buterin2014}은 스마트 계약을 통해 거버넌스를 구현하지만, 영향력이 자본에 비례하는 금권정치적 역학을 재현한다. 이차 투표(quadratic voting)~\cite{buterin2019}는 이를 부분적으로 완화하지만 실질적 채택을 제한한다. DAO는 또한 근본적인 \emph{오라클 문제}에 직면한다. 스마트 계약은 신뢰할 수 있는 외부 출처 없이는 현실 세계의 상태를 안정적으로 검증할 수 없으며, 이 문제는 블록체인 아키텍처 내에서 일반적인 해법이 없다.

기존의 어떤 시스템도 탈중앙화, 검열 저항성, 가명성, 검증 가능한 진입 비용, 그리고 창발적 합의를 결합하지 못한다.

\subsection{기여}

본 논문은 다음과 같은 기여를 한다.
\begin{enumerate}[nosep]
\item DID 프레임워크를 확장하여 양방향 평판 교환을 지원하는 탈중앙화 프로토콜인 \textbf{평판 소셜 네트워크(RSN)}의 정의.
\item 일관된 해싱~\cite{karger1997}에 기반한 \textbf{비결정론적 검증자 선택 프로토콜}.
\item 세 가지 복합적 비용 채널을 통해 네트워크 합의로부터의 거리에 따라 주장에 가격을 매기는 \textbf{급진주의 비용화 원칙(Expensive Radicalism Principle)}.
\item 거짓 보증이 정당한 수수료보다 파국적으로 더 큰 비용을 치르게 하는 비대칭적 유인을 지닌 \textbf{평판 권위(Reputable Authority) 모델}.
\item 병렬 재정 인프라를 통한 \textbf{점진적 이행 경로}.
\end{enumerate}

% ============================================================
% 2. DESIGN PRINCIPLES
% ============================================================
\section{설계 원칙}

RSN 아키텍처는 다섯 가지 타협 불가능한 설계 원칙에 의해 제약된다.

\textbf{연속성과 진화.} 시스템은 무기한의 이행 기간 동안 기존 기관과 공존한다. 이행은 모든 단계에서 되돌릴 수 있어야 한다.

\textbf{자발성과 책임.} 참여는 자발적이지만, 자발성은 책임과 결합된다. 제도적 기능에 대한 통제권을 떠맡는 참여자는 동시에 그에 수반되는 의무도 떠맡는다.

\textbf{다양성과 문화적 중립성.} 합의는 시스템 설계자가 부과한 사전 규칙이 아니라 양방향 상호작용에서 창발한다.

\textbf{회복력과 검열 저항성.} 네트워크를 검열하는 비용은 그것을 용인하는 비용을 초과해야 한다. 오직 완전한 전체주의만이---파멸적인 정치적, 경제적 비용을 치르고---시스템을 억압할 수 있다.

\textbf{합의와 집행.} 시스템은 옹졸함, 앙심, 응보적 만족을 향한 인간의 경향을 하중을 지탱하는 핵심 요소로 통합한다. 부정행위는 금지되는 것이 아니라 가격이 매겨진다.

% ============================================================
% 3. SYSTEM OVERVIEW
% ============================================================
\section{시스템 개요}

\subsection{평판 소셜 네트워크}

RSN은 참여자들이 현실 세계의 행동에 관한 검증 가능한 정보를 교환하는 탈중앙화된 P2P 통신 계층이다. 그 본질적 속성은 다음과 같다. 탈중앙화되고 검열할 수 없는 인프라, DID를 통한 가명 참여, 비대칭적 비용 구조(읽기는 저렴하고 쓰기는 비쌈), 암호학적 검증 가능성, 그리고 \textbf{표준 지원}이다. 표준 지원은 네트워크를 통해 전파되는 정보, 권위가 제공하는 서비스(주장 구성, 보증, 증인 서비스), 그리고 상대방 평판을 평가하는 규칙과 (선언된 행동과 실제 행동 간의) 정책 불일치 위험을 평가하는 규칙을 포함한 정책 형식에 대한 기계 판독 가능한 형식을 뜻한다.

\subsection{아키텍처 구성 요소}

RSN은 네 가지 주요 구성 요소로 이루어진다(그림~\ref{fig:architecture}).
\begin{enumerate}[nosep]
\item \textbf{DID 계층.} 선언된 규칙, 평판 메타데이터, 네트워크 라우팅을 갖춘 참여자 신원.
\item \textbf{주장 프로토콜(Claim Protocol).} 현실 세계 사건에 관한 단언을 게시하기 위한 구조화된 형식.
\item \textbf{검증 네트워크.} 일관된 해싱을 사용하여 검증자를 배정하는 탈중앙화 프로토콜.
\item \textbf{평판 집계 계층.} 사람이 읽을 수 있는 평판 요약을 생성하는 서비스.
\end{enumerate}

\begin{figure}[ht]
\centering
\begin{tikzpicture}[
    layer/.style={draw, minimum height=0.7cm, align=center, font=\scriptsize, fill=black!8},
    arr/.style={<->, thick, gray!60}
]
\node[layer, minimum width=5cm] (agg) at (0,2.8) {\textbf{집계}\\해석\ $\cdot$ 위험};
\node[layer, minimum width=2.2cm] (claim) at (-1.55,1.4) {\textbf{주장}\\구성};
\node[layer, minimum width=2.2cm] (verif) at (1.55,1.4) {\textbf{검증}\\선택};
\node[layer, minimum width=5cm] (did) at (0,0) {\textbf{DID 계층}\\신원 $\cdot$ 규칙 $\cdot$ 점수};

\draw[arr] (agg.south west) -- (claim.north);
\draw[arr] (agg.south east) -- (verif.north);
\draw[arr] (claim.east) -- (verif.west);
\draw[arr] (claim.south) -- (did.north west);
\draw[arr] (verif.south) -- (did.north east);
\end{tikzpicture}
\caption{RSN 4계층 아키텍처.}
\label{fig:architecture}
\end{figure}

\subsection{참여자 역할}

검증 거래에는 최대 여섯 개의 구별되는 DID 역할이 관여한다(그림~\ref{fig:roles}). \textbf{발행자}($DID_I$, 주장을 생성하고 제출), \textbf{대상}($DID_S$, 주장이 관한 DID---발행자와 일치할 수 있음), \textbf{권위}($DID_A$, 수수료를 받고 주장을 보증하며 증인 서비스도 제공할 수 있음---\emph{선택 사항}), \textbf{증인}($DID_{W_{1..n}}$, 챌린지 코드를 통해 검증자의 정직성을 익명으로 감시---\emph{선택 사항}), \textbf{검증자}($DID_V$, 주장을 게시하도록 알고리즘으로 선택됨), 그리고 \textbf{RSN}(검증된 주장이 게시되는 네트워크)이다. 모든 역할은 다른 DID에 위임될 수 있다(7.4절). 권위는 통상 보증과 증인 서비스를 묶어서 제공하지만, 두 기능은 논리적으로 독립적이다.

\begin{figure}[ht]
\centering
\begin{tikzpicture}[
    role/.style={draw, rounded corners, minimum width=1.1cm, minimum height=0.45cm, align=center, font=\scriptsize, fill=black!8},
    opt/.style={draw, rounded corners, minimum width=1.1cm, minimum height=0.45cm, align=center, font=\scriptsize, fill=black!8, dashed},
    arr/.style={-{Stealth[length=3pt]}, thick},
    oarr/.style={-{Stealth[length=3pt]}, thick, dashed, gray},
    lbl/.style={font=\tiny, fill=white, inner sep=1pt}
]
\node[role] (I) at (0,2.4) {\textbf{발행자}};
\node[role] (S) at (-2.5,2.4) {\textbf{대상}};
\node[opt] (A) at (2.5,1.2) {\textbf{권위}};
\node[opt] (W) at (-2.5,0) {\textbf{증인}};
\node[role] (V) at (2.5,-0.6) {\textbf{검증자}};
\node[role, fill=black!5, draw=gray] (N) at (0,-1.8) {RSN};

\draw[arr] (I) -- node[lbl] {관함} (S);
\draw[oarr] (I) -- node[lbl,sloped] {보증} (A);
\draw[oarr] (I) -- node[lbl,sloped] {챌린지} (W);
\draw[arr] (I) -- node[lbl,sloped] {요청} (V);
\draw[oarr] (W) -- node[lbl] {감시} (V);
\draw[arr] (V) -- node[lbl] {게시} (N);
\end{tikzpicture}
\caption{검증 거래에서의 DID 역할. 점선 = 선택 사항(권위, 증인).}
\label{fig:roles}
\end{figure}

\subsection{부패 불가능성: 네 가지 힘}

RSN의 부패 불가능성은 단일 메커니즘에서 비롯되는 것이 아니라 동시에 작용하는 네 가지 힘에서 비롯된다. 어느 것도 그 자체로 충분하지 않다. 오직 그것들의 결합만이 부패 시도를 경제적으로 비합리적으로 만든다.

\textbf{예측 불가능한 표적.} 각 주장의 검증자는 일관된 해싱을 통해 비결정론적으로 선택된다(5.3절). 검증자의 신원은 발행자의 선택이 아니라 주장의 내용과 이전 반복의 함수이므로, 공격자는 매수를 위해 특정 검증자를 미리 겨냥할 수 없다.

\textbf{평판 레버리지---상승은 느리고, 하락은 빠르다.} 평판은 지속적인 일관된 행동을 통해 점진적으로만 얻을 수 있다. 그러나 단 한 번의 사기성 보증으로 파국적으로 잃을 수 있다(6.2절). 이 비대칭은 수년간 축적한 평판이 단 한 번의 부정직한 보증으로 파괴될 수 있음을 뜻한다. 최적의 전략은 여전히 의심스러운 주장을 거부하는 것이다.

\textbf{공개적 약속.} 모든 DID 보유자는 자신의 정책---따르기로 약속한 기계 판독 가능한 규칙의 집합---을 공개적으로 선언한다. 선언과 실제 행동 사이의 괴리는 탐지 가능하며 근본적인 위반보다 더 심각한 평판 위반으로 가격이 매겨진다(7.3절). 따라서 위선은 직접적인 부정직함보다 더 큰 비용을 치른다.

\textbf{메시(mesh) 공격 비용 비대칭.} 네트워크를 검열하거나 불안정하게 만드는 비용은 네트워크의 규모와 밀도에 따라 비선형적으로 증가하는 반면, 그것을 용인하는 비용은 일정하게 유지된다. 검열이 수지를 맞추려면 중앙집권적 행위자가 그것을 네트워크 전체에 적용해야 하는데, 이는 어떤 상상 가능한 이익과도 비례하지 않는 조치를 요구한다(10절).

% ============================================================
% 4. DECENTRALIZED IDENTITY MODEL
% ============================================================
\section{탈중앙화 신원 모델}

\subsection{특수 속성을 지닌 DID}

RSN은 W3C DID Core 명세~\cite{did2022}를 다음으로 확장한다. \textbf{선언된 규칙}(기계 판독 가능한 행동 약속), \textbf{평판 메타데이터}(집계된 행동 점수), 그리고 \textbf{네트워크 라우팅}(안정적인 링 위치와 분리된 가변적 어니언 게이트웨이).

\subsection{주장 생애주기}

주장은 여섯 단계를 거친다(그림~\ref{fig:lifecycle}). 구성, 선택적 권위 보증, 일관된 해싱을 통한 검증자 선택, 검증 결정(수락 또는 반복과 함께 거부), 게시, 그리고 모든 당사자에 대한 평판 갱신이다.

\begin{figure}[ht]
\centering
\begin{tikzpicture}[
    stp/.style={draw, rounded corners=2pt, minimum width=1.3cm, minimum height=0.45cm, align=center, font=\tiny, fill=black!5},
    arr/.style={-{Stealth[length=3pt]}, thick},
    lbl/.style={font=\tiny, fill=white, inner sep=1pt}
]
\node[stp] (comp) at (0,0) {주장\\구성};
\node[stp, dashed] (auth) at (2.0,0) {권위\\보증};
\node[stp] (hash) at (4.0,0) {해시 \&\\링 매핑};
\node[stp] (ver) at (2.0,-1.6) {검증자\\확인};
\node[stp, double] (pub) at (0,-1.6) {게시됨};
\node[stp] (rej) at (4.0,-1.6) {거부됨\\비용 $\uparrow$};

\draw[arr] (comp) -- (auth);
\draw[arr] (auth) -- (hash);
\draw[arr] (hash) -- (ver);
\draw[arr] (ver) -- node[lbl] {\checkmark} (pub);
\draw[arr] (ver) -- node[lbl] {$\times$} (rej);
\draw[arr] (rej) -- ++(0,0.8) -| (hash);
\end{tikzpicture}
\caption{반복 루프를 지닌 주장 생애주기.}
\label{fig:lifecycle}
\end{figure}

\subsection{W3C DID Core와의 관계}

RSN의 신원 모델은 W3C DID Core 명세~\cite{did2022}의 진부분집합의 상위 집합이다. 모든 RSN DID는 유효한 W3C DID다. RSN 고유의 확장은 전용 DID 방식 명세에 정의된 DID 문서 속성으로 인코딩되며, ION~\cite{buchner2021}과 Ceramic~\cite{ceramic2023} 같은 기존 인프라와 호환된다.

% ============================================================
% 5. VERIFICATION AND PROPAGATION
% ============================================================
\section{정보 검증과 전파}

\subsection{정보 진입 비용}

RSN의 핵심 경제 원칙은 읽기와 쓰기 사이의 비대칭이다. 평판 데이터를 읽는 것은 DID 네트워크의 일부이고 자신의 평판에 상응하는 접근 권한을 가진 참여자에게 한계 비용이 0에 가깝다---신규 참여자는 점진적으로 더 넓은 접근 권한을 얻어야 한다(반(反)무임승차 참조). 쓰기---일회성 기술적 행위가 아니라 평판을 네트워크에 지속적으로 구현하는 것으로 이해되는---는 세 차원에 걸친 검증 가능한 누적 투자를 요구한다. \textbf{시간}(일관된 행동, 이행된 약속, 가꾼 관계에 쓰인 인생의 세월), \textbf{에너지}(정직한 거래, 협력 관계에 대한 배려, 장기적 신뢰성에 투자한 노력), 그리고 \textbf{돈}(자신의 이름에 대한 투자---전문성 개발, 자기 일의 품질, 초래한 피해에 대한 배상)이다. 평판은 즉시 살 수 없다---그것은 시스템이 그저 가시화하고 맥락을 넘어 이전 가능하게 만들 뿐인, 평생에 걸친 축적의 결과다. 프로토콜의 일회성 기술 비용(암호학적 서명, 검증자 수수료)은 이 근본적 투자에 비하면 무시할 만하다.

\subsection{소셜 그래프와 접촉 깊이}

RSN은 근본적으로 소셜 네트워크다. 각 DID는 연결을 허용하는 접촉자(다른 DID)를 추가한다. 이 접촉자들은 자신의 접촉자를 가지며, 그런 식으로 이어진다. 알고리즘적 검증자 탐색은 연결된 접촉자의 설정 가능한 깊이 $h$ 이내에서 작동한다(예: $h=3$: 자신의 직접 접촉자, 그들의 접촉자, 그리고 접촉자의 접촉자). 이 아키텍처는 단일 글로벌 블록체인을 요구하지 않는다---그것은 자연스럽게 다른 커뮤니티와 겹치는 커뮤니티/클러스터의 창발을 선호한다. 모든 DID 참여자는 게시된 \textbf{정책}---들어오는 요청을 어떻게 평가할지 지배하는 기계 판독 가능한 규칙의 집합---을 가져야 한다. 정책 위반은 부정적 아티팩트로 네트워크에 기록된다.

\subsection{알고리즘적 검증자 선택}

검증 프로토콜은 일관된 해싱~\cite{karger1997}을 사용한다(그림~\ref{fig:verifier}). 검증은 유료 서비스다. 검증자는 수수료를 받고 운영하며, 이는 경쟁이 가격, 속도, 신뢰성을 견인하는 시장을 만든다.

\textbf{1단계.} 주장 문서는 이전 반복의 해시 결과(또는 첫 반복에서는 $\varnothing$)와 연결되어 해싱된다.
\begin{equation}
\text{pos} = f_h(\text{claim} \| \text{prev\_hash})
\label{eq:hash}
\end{equation}

알고리즘은 단지 이전 반복의 해시가 아니라 모든 이전 요청과 응답의 \emph{완전한 경로}를 포함해야 한다. 각 검증자의 전체 응답(수락, 거부, 제시된 가격, 사유)은 커져 가는 문서에 추가되며, 각 단계에서 암호학적 서명을 통해 검증 가능하다.

\textbf{2단계.} 해시는 일관된 해시 링 위의 $N$개 위치에 매핑되며, 이는 균등하게 분포된다(예: $N=4$의 경우 $+0^{\circ}, +90^{\circ}, +180^{\circ}, +270^{\circ}$). 각 위치에서 시계 방향 탐색이 가장 가까운 DID를 검증자 후보로 선택한다(그림~\ref{fig:multipoint}). $N$은 현재 활성 DID의 비율, 하루 중 시각, 또는 과거 네트워크 응답성에 따라 달라질 수 있는 가변 매개변수다.

\textbf{3단계.} 검증자는 (평판을 걸고 게시하며) 수락하거나, (거부 해시와 함께 1단계로 돌아가) 거부한다. 온라인 상태를 선언했음에도 노드가 응답하지 않을 때, 다음 요청은 이전 $N$개 검증자 후보 모두의 응답을 포함해야 한다.

\textbf{반(反)무임승차.} 대상 DID는 평판이 적은 신규 DID의 조회에 대해 수수료나 접근 제한을 설정할 수 있다. 이 점진적 메커니즘(이분법이 아님)은 신규 참여자가 타인의 데이터를 자유롭게 소비하기 전에 진정한 활동을 통해 평판을 쌓도록 강제한다.

\begin{figure}[ht]
\centering
\begin{tikzpicture}[
    box/.style={draw, rounded corners=2pt, minimum width=1.3cm, minimum height=0.45cm, align=center, font=\scriptsize, fill=black!5},
    dec/.style={draw, diamond, aspect=2.2, minimum width=0.9cm, align=center, font=\scriptsize, fill=black!5},
    arr/.style={-{Stealth[length=3pt]}, thick},
    lbl/.style={font=\tiny, fill=white, inner sep=1pt}
]
\node[box] (iss) at (0,0) {발행자};
\node[box] (hash) at (0,-1.1) {$f_h$(claim $\|$\\prev\_hash)};
\node[box] (ring) at (0,-2.2) {링 매핑\\시계 방향};
\node[box, dashed] (spam) at (0,-3.2) {반스팸\\보증금};
\node[dec] (check) at (0,-4.4) {정책\\확인};
\node[box] (rej) at (2,-5.6) {다음\\반복};
\node[box] (fee) at (0,-5.6) {최종 수수료 =\\$f$(데이터, 평판, 정책)};
\node[box, double] (pub) at (0,-6.8) {게시됨};

\draw[arr] (iss) -- (hash);
\draw[arr] (hash) -- (ring);
\draw[arr] (ring) -- (spam);
\draw[arr] (spam) -- (check);
\draw[arr] (check) -- node[lbl] {$\times$} (rej);
\draw[arr] (check) -- node[lbl] {\checkmark} (fee);
\draw[arr] (fee) -- (pub);
\draw[arr] (rej.east) -- ++(0.5,0) |- (hash.east);
\end{tikzpicture}
\caption{검증자 선택 프로토콜. 반스팸 보증금은 선택 사항이며 환불 가능할 수 있다. 최종 수수료는 수락 시에만 지불된다.}
\label{fig:verifier}
\end{figure}

각 거부는 검증자의 전체 응답(사유와 조건 포함)을 주장 문서에 추가하여 그 내용을 확장한다. 확장된 문서로부터 새로운 해시가 비결정론적으로 계산되어 다른 링 위치에 매핑되고 다른 검증자를 선택한다. 전체 경로의 문서화는 필수다---발행자는 다음 검증자를 예측하거나 영향을 미칠 수 없다. 검증자가 호환되는 선언 정책에도 불구하고 요청을 무시하면, (증인이 있는 경우) 증인이 이를 기록하고, 발행자는 최종 게시 시 증인 증거를 첨부할 수 있다.

\begin{figure}[ht]
\centering
\begin{tikzpicture}[scale=0.65]
% Ring
\draw[thick] (0,0) circle (2.2cm);
% DID nodes on ring
\foreach \a in {10,55,80,120,150,195,230,260,305,340} {
  \fill[black!40] (\a:2.2) circle (1.5pt);
}
% Hash offset arrow pointing to initial position
\draw[-{Stealth[length=3pt]}, thick, black!60] (0,0) -- node[font=\tiny,above,sloped,pos=0.6] {$f_h$} (37:2.2);
\fill[black] (37:2.2) circle (2pt);
\node[font=\tiny,anchor=south west] at (37:2.35) {$h_0$};
% N=4 points offset from h0 by +0, +90, +180, +270
\foreach \offset/\l in {0/P1,90/P2,180/P3,270/P4} {
  \pgfmathsetmacro{\ang}{37+\offset}
  \node[fill=black, circle, inner sep=1.5pt] (p\offset) at (\ang:2.2) {};
  \node[font=\tiny,font=\bfseries] at (\ang:2.75) {\l};
  % clockwise search arc
  \draw[-{Stealth[length=2.5pt]}, thick, gray] (\ang:2.2) arc (\ang:\ang+30:2.2);
}
\node[font=\scriptsize, align=center] at (0,0) {해시\\링};
\node[font=\tiny, align=center] at (0,-3.2) {$h_0 = f_h(\text{claim})$, 이후 $+0^{\circ},+90^{\circ},+180^{\circ},+270^{\circ}$\\각 지점 $\rightarrow$ 시계 방향 탐색};
\end{tikzpicture}
\caption{다중 지점 선택($N{=}4$). 해시 $h_0$가 오프셋을 설정하고, $h_0$에서 균등하게 분포된 4개 지점.}
\label{fig:multipoint}
\end{figure}

\subsection{증인 프로토콜}

\textbf{증인} 역할은 암호학적 챌린지 코드 프로토콜을 통해 검증자의 정직성에 대한 익명의 책임성을 제공한다.

\textbf{W1단계.} 요청자는 검증 요청에 서명하고 이를 $N_w$명의 증인---요청자의 소셜 네트워크의 접촉자, 또는 수수료를 받고 운영하는 전문 증인 서비스 제공자---에게 보낸다.

\textbf{W2단계.} 각 증인 $w_i$는 서명된 전체 요청에 서명하고, 원본 서명 $\sigma_i$를 보관하며, 챌린지 코드 $c_i = h(\sigma_i)$를 계산한다. 챌린지 코드는 요청에 추가된다.

\textbf{W3단계.} 이제 $N_w$개의 챌린지 코드를 지닌 요청이 검증자에게 전송된다. 검증자는 코드를 관찰하지만 누가 증인인지, 또는 코드가 실제 서명에 대응하는지 \emph{판별할 수 없다}.

\textbf{W4단계(정직한 검증자).} 검증자가 선언된 정책에 따라 응답하면, 챌린지 코드는 불투명하게 남는다. 증인은 결코 드러나지 않는다. 추가 데이터는 게시되지 않는다.

\textbf{W5단계(정책 위반).} 검증자가 정책을 위반하면(요청 무시, 일관되지 않은 응답), 요청자는 증인의 원본 서명 $\sigma_i$를 보유하고 있다. 이는 동반 데이터로 게시될 수 있다. 누구든 $c_i = h(\sigma_i)$를 검증하여 증인이 있었음을 증명할 수 있다. 검증자는 이미 응답에서 챌린지 코드에 서명했으므로 요청자는 독립적으로 게시할 수 있다.

\textbf{블러프(bluff) 속성.} 요청자는 일부 또는 모든 챌린지 코드를 무작위 값으로 대체할 수 있다. 검증자는 진짜 코드(고평판 권위가 뒷받침하는)를 잡음과 구별할 수 없다. 이 \emph{불확실성}이 집행 메커니즘이다. 모든 요청은 존경받는 권위가 익명으로 지켜보고 있을 위험을 수반한다. 이 압박을 유지하는 비용은 0에 가깝지만(무작위 숫자 생성), 검증자에게 부정직함이 초래할 잠재적 비용은 파국적이다. 실제 증인이 없더라도 정직한 행동이 장려된다.

\textbf{증인으로서의 권위.} 주장을 보증하는 권위는 증인 서비스를 묶을 수 있다. ``나는 당신의 주장을 보증하며 검증 동안 익명의 증인으로 활동한다.'' 이는 자연스러운 사업 모델이다---권위는 이미 맥락을 가지고 있다. 그러나 두 역할은 논리적으로 독립적이다.

\subsection{급진주의 비용화 원칙}

세 가지 비용 채널이 동시에 복합된다(그림~\ref{fig:radicalism}).

\textbf{채널 1: 반복 비용.} 각 거부는 시간과 자원을 소비하는 새로운 반복을 강제한다. 선형 증가.

\textbf{채널 2: 문서 팽창.} 각 반복은 검증자의 전체 응답을 주장 문서에 추가한다. 증가는 선형이지만 기저 오프셋이 있다. 초기 페이로드(주장 내용, 권위 보증, 암호학적 서명)는 이미 무시할 수 없는 크기 $B_0$를 가진다. $k$번 반복 후 전체 크기는 $S(k) = B_0 + k \cdot \Delta$이며, 여기서 $\Delta$는 평균 응답 크기다.

\textbf{채널 3: 검증자 위험 프리미엄.} 위험은 주관적이다. 검증자는 요청하는 DID의 선언된 정책이 검증자 자신의 상업적, 사회적, 또는 정치적 이해와 충돌하는지 평가한다. 프리미엄은 검증자의 ``이상''으로부터의 인지된 거리에 따라 지수적으로 상승할 수 있다. $P(d) = \alpha \cdot e^{\beta d}$, 여기서 $d$는 검증자의 주관적 거리 척도다. 개인적 금지선(no-go boundary)을 넘어서면 검증자는 전면적으로 거부할 수 있다.

\begin{figure}[ht]
\centering
\begin{tikzpicture}
\begin{axis}[
    width=\columnwidth,
    height=4.5cm,
    xlabel={\footnotesize 급진성},
    ylabel={\footnotesize 상대 비용 (로그)},
    ymode=log,
    xmin=0, xmax=100,
    ymin=1, ymax=10000,
    grid=major,
    grid style={gray!20},
    legend style={at={(0.03,0.97)},anchor=north west,font=\tiny,draw=gray!50},
    tick label style={font=\tiny},
    label style={font=\footnotesize},
]
\addplot[black, thick, domain=0:100, samples=50] {1 + 0.3*x};
\addlegendentry{반복 비용 (선형)}
\addplot[black, thick, dashed, domain=0:100, samples=50] {1 + 0.15*x};
\addlegendentry{문서 팽창 (선형)}
\addplot[black, thick, dotted, line width=1.2pt, domain=0:100, samples=100] {exp(0.08*x)};
\addlegendentry{위험 프리미엄 (지수)}
\end{axis}
\end{tikzpicture}
\caption{세 채널에 걸친 비용 상승. 위험 프리미엄이 높은 급진성에서 지배적이다.}
\label{fig:radicalism}
\end{figure}

결정적으로, 이것은 검열이 아니다. 어떤 주장도 금지되지 않는다. 시스템은 위험에 가격을 매긴다(표~\ref{tab:costs}).

\begin{table}[ht]
\centering
\caption{주장 유형별 비용 비교.}
\label{tab:costs}
\footnotesize
\begin{tabular}{@{}lcccc@{}}
\toprule
\textbf{유형} & \textbf{반복} & \textbf{문서} & \textbf{수수료} & \textbf{총계} \\
\midrule
신뢰성 있음 & $k_{\min}$ & $B_0$ & $P_{\min}$ & 낮음 \\
논쟁적 & $k_{\text{med}}$ & $B_0{+}k\Delta$ & $\alpha e^{\beta d}$ & 중간 \\
급진적 & $k_{\max}$ & $\gg B_0$ & $\gg P_{\min}$ & 매우 높음 \\
사기성 & $\to\infty$ & --- & --- & 게시 불가 \\
\bottomrule
\end{tabular}
\end{table}

% ============================================================
% 6. REPUTATION AND RISK
% ============================================================
\section{평판과 위험 평가}

\subsection{평판 축적 모델}

평판은 세 가지 속성을 보인다. \textbf{느린 축적}(일관된 행동을 통한 점진적 성장), \textbf{빠른 파괴}(단 한 번의 사기가 점수를 파국적으로 떨어뜨릴 수 있음), 그리고 \textbf{개별적 평가}(각 소비 당사자가 자신만의 집계 함수를 적용할 수 있음)다.

\subsection{평판 권위}

평판 권위는 주장을 검토하고, 만족하면 이를 공동 서명한다---자신의 평판을 건다(그림~\ref{fig:authority}). 비대칭은 의도된 설계다. 평판을 얻는 것은 느리지만(정직한 보증마다 점진적 $+\delta$), 잃는 것은 파국적이다(거짓 보증마다 $-\Delta \gg \delta$; 예시로 $+2\%$ 대 $-70\%$). 최적의 전략은 언제나 의심스러운 주장을 거부하는 것이다. $\delta$와 $\Delta$의 구체적 값은 시스템 매개변수다.

\begin{figure}[ht]
\centering
\begin{tikzpicture}[
    box/.style={draw, rounded corners=2pt, minimum width=2cm, minimum height=0.6cm, align=center, font=\scriptsize, fill=black!5},
    arr/.style={-{Stealth[length=4pt]}, thick}
]
\node[box] (exam) at (0,0) {권위가 검토\\평판: $R$};
\node[box] (true) at (-2,-1.3) {참: $R{\to}R{+}\delta$\\고객 증가};
\node[box] (false) at (2,-1.3) {거짓: $R{\to}R{-}\Delta$\\고객 이탈};
\node[box] (insight) at (0,-2.6) {획득: 느림 ($+\delta$)\\상실: 즉각 ($-\Delta$)};

\draw[arr] (exam) -- node[left,font=\tiny] {진실} (true);
\draw[arr] (exam) -- node[right,font=\tiny] {거짓} (false);
\draw[arr] (true) -- (insight);
\draw[arr] (false) -- (insight);
\end{tikzpicture}
\caption{평판 권위---비대칭적 유인.}
\label{fig:authority}
\end{figure}

\subsection{다차원 평판}

평판은 스칼라가 아니라 여러 차원에 걸친 벡터다. 재정적 신뢰성, 개인적 성실성, 전문적 역량, 시민적 참여 등(그림~\ref{fig:reputation-radar}). 서로 다른 평가 제공자는 맥락에 따라 이 벡터를 다르게 투영한다---대출 기관은 재정적 신뢰성을, 고용주는 전문적 역량을, 이웃은 시민적 행동을 우선시한다. 단 하나의 ``올바른'' 평판 점수는 없다. 오직 관점만이 있을 뿐이다.

\begin{figure}[ht]
\centering
\begin{tikzpicture}[scale=0.55]
\foreach \a/\l in {90/재정,162/성실,234/전문,306/시민,18/사회} {
  \draw[gray!40] (0,0) -- (\a:2.5);
  \node[font=\tiny, align=center] at (\a:3.1) {\l};
  \foreach \r in {0.8,1.6,2.4} {
    \fill[black!15] (\a:\r) circle (0.5pt);
  }
}
\draw[thick, fill=black!10] (90:2.0) -- (162:1.2) -- (234:1.8) -- (306:0.9) -- (18:1.5) -- cycle;
\draw[thick, dashed] (90:1.4) -- (162:2.1) -- (234:0.8) -- (306:2.0) -- (18:1.1) -- cycle;
\node[font=\tiny] at (0,-3.3) {실선: DID A \quad 점선: DID B};
\end{tikzpicture}
\caption{다차원 평판 벡터. 서로 다른 DID는 차원에 걸쳐 서로 다른 프로필을 보인다.}
\label{fig:reputation-radar}
\end{figure}

\subsection{위험 평가의 민주화}

RSN은 체계적 위험 평가를 민주화한다---현재 금융 기관에 집중되어 있지만 은행업을 훨씬 넘어 적용 가능한 역량이다. 공급자 신뢰성을 평가하는 산업 대기업, 행동 위험에 가격을 매기는 보험사, 인수합병을 위한 실사, 제조업에서의 장비 수명 추정---상대방 위험 평가가 필요한 곳이면 어디서든 RSN은 탈중앙화되고 시장 주도적인 대안을 제공한다. 해석 서비스 시장은 원시 다차원 평판 데이터를 실행 가능한 요약으로 번역하여, 어떤 참여자든 한계 비용으로 상대방 평가에 접근할 수 있게 한다.

% ============================================================
% 7. CONSENSUS AND DISPUTE RESOLUTION
% ============================================================
\section{합의와 분쟁 해결}

\subsection{탈중앙화 정의 모델}

RSN은 정의를 양방향 협상 문제로 재구성한다. 영구적이고 검증 가능한 평판 훼손의 위협은 피해를 입은 당사자가 감당할 수 없는 법적 절차보다 더 효과적인 억지력이다.

\subsection{창발적 합의}

각 참여자는 개인적 만족 문턱을 유지한다. 네트워크의 집합적 합의는 수천 건의 양방향 협상의 통계적 평균으로 창발한다(그림~\ref{fig:consensus}). 합의를 크게 상회하는 대응(야만적)은 게시하기에 비싸다. 합의에 근접한 대응(비례적)은 저렴하다. 합의는 규칙이 아니다---그것은 가격 신호다.

\begin{figure}[ht]
\centering
\begin{tikzpicture}[
    person/.style={draw, rounded corners=2pt, minimum width=1.5cm, minimum height=0.5cm, align=center, font=\scriptsize, fill=black!5},
    arr/.style={-{Stealth[length=4pt]}, thick}
]
\node[person] (a) at (-2,0) {엄격\\(높음)};
\node[person] (b) at (0,0) {실용적\\(중간)};
\node[person] (c) at (2,0) {관대\\(낮음)};
\node[person, fill=black!10, minimum width=3cm] (cons) at (0,-1.2) {네트워크 합의\\= 통계 평균};
\node[person] (exp) at (-2,-2.4) {야만적\\비쌈};
\node[person, double] (prop) at (0,-2.4) {비례적\\저렴};
\node[person] (neg) at (2,-2.4) {태만\\비쌈};

\draw[arr] (a) -- (cons);
\draw[arr] (b) -- (cons);
\draw[arr] (c) -- (cons);
\draw[arr] (cons) -- (exp);
\draw[arr] (cons) -- (prop);
\draw[arr] (cons) -- (neg);
\end{tikzpicture}
\caption{개인적 만족 문턱으로부터의 창발적 합의.}
\label{fig:consensus}
\end{figure}

\subsection{처벌 보정}

각 DID 보유자는 특정 부정행위 범주에 대한 대응을 공개적으로 선언한다. 불균형하게 가혹하거나 관대한 선언은 부정적 평판 신호를 끌어들인다. 비례적 선언은 마찰 없이 수용된다---자기보정적 처벌 시스템이다.

합의 선언 자체도 위선 탐지의 대상이다. 합의 입장에 선언적으로 약속하면서 일관되지 않게 행동하는 것은 의도적 기만을 드러내므로, 근본적 일탈보다 더 심각한 평판 위반을 구성한다.

\subsection{위임}

DID 내의 모든 활동은---한 가지 예외를 제외하고---다른 DID에 위임될 수 있다. \textbf{대상 역할---주장이 그 행동에 관한 DID---은 위임될 수 없다. 그것은 주장이 지칭하는 신원 보유자에게 양도 불가능하게 결속된다.} 다른 능동적 역할---발행자, 권위, 증인, 검증자---은 검증, 주장 제출, 합의 참여, 또는 분쟁 해결을 위해 지정된 대리 DID에 이전될 수 있다. 위임은 언제든 철회 가능하다. 이는 보유자가 궁극적 통제권과 책임을 유지하면서 전문화(예: 전문 검증 서비스)를 가능하게 한다. 위임은 시스템 전반에 적용되며 확장성에 필수적이다.

\subsection{개인의 도덕에서 창발적 사회 계약으로}

각 DID의 선언된 정책은 개인 도덕의 형식화를 나타낸다. 즉 보유자가 무엇을 수용 가능하다고 여기는지, 특정 행동에 어떻게 대응하는지, 상대방에게 무엇을 요구하는지다. 이 정책들은 시스템 설계자가 부과하는 것이 아니다---각 참여자가 선택하고 기계 판독 가능한 형식으로 표현한다.

이 형식화는 세 단계의 창발을 가능하게 한다. 첫째, 개인 도덕은 \textbf{정책}---DID가 따르기로 약속한 선언된 규칙, 문턱, 대응 함수의 집합---으로 인코딩된다. 둘째, 네트워크 전반의 모든 정책의 집합은 \textbf{창발적 윤리}---어떤 단일 참여자도 설계하지 않았지만 수천 개의 개별 도덕적 입장의 상호작용에서 발생하는, 통계적으로 관찰 가능한 규범~\cite{durkheim1893}---를 산출한다. 셋째, 양방향 가격 신호를 통해 집행되는 공유된 행동 규범으로 표현된 이 창발적 윤리는 \textbf{측정 가능한 사회 계약}을 구성한다---아무도 서명하지 않은 이론적 구성물~\cite{rawls1971}이 아니라, 실제 행동으로 뒷받침되는 경험적으로 관찰 가능한 규칙의 집합이다.

이 과정은 도덕 기반 이론~\cite{haidt2012}의 발견을 반영한다. 인간의 도덕은 직관적이고, 다원적이며, 문화적으로 가변적이다. RSN은 입력으로 도덕적 합의를 요구하지 않는다. 그것은 출력으로 행동적 합의를 산출한다. 그 결과로 나온 사회 계약은 정적이지 않다---참여자가 가입하고, 탈퇴하고, 네트워크 피드백에 반응하여 정책을 조정함에 따라 진화한다. 고전적 사회 계약 이론과 달리, 이 계약은 철회 가능하고, 관찰 가능하며, 주권자 없이도 집행 가능하다.

% ============================================================
% 8. ECONOMIC MODEL
% ============================================================
\section{경제 모델}

앞선 절들은 도구---RSN의 검증 메커니즘, 비용 구조, 창발적 합의---를 설명했다. 이 절은 이 도구를 재정 및 거버넌스 이행에 적용하는 방법론을 설명한다. 도구는 범용적이며, 아래의 방법론은 하나의 가능한 응용이다.

\subsection{유인 구조}

자본으로서의 평판은 정직한 행동에 대한 직접적 유인을 만든다. 정상적인 운영에서 참여자는 일상적 거래와 이행된 약속을 통해 자연스럽게 평판을 쌓는다. 주된 지출은 \emph{부정적} 주장---타인이 저지른 피해를 기록하는 것---에 쓰인다. 이 비대칭은 유용하고 신뢰할 수 있는 행동을 장려한다. 정직한 것은 추가 비용이 들지 않는 반면, 해로운 것은 타인이 보존하는 데 기꺼이 비용을 치를 문서화된 흔적을 만든다. 위선---규칙을 선언하고 따르지 않는 것---은 가장 비싼 실패 양태다. 그것이 의도적 기만을 드러내기 때문이다.

\subsection{병렬 재정 시스템}

세 가지 구성 요소가 경쟁 압력을 가능하게 한다. (1)~\textbf{단순세}---예외 없는 단일 비례 세율로, 초기에는 기존 실효 세율보다 근소하게 낮게. (2)~\textbf{전자 지출 등록(ESR)}---체코의 EET 모델을 뒤집어 시민이 국가 지출을 감시하도록 함. (3)~\textbf{시민 지정 세금 배분}---첫해에 몇 퍼센트로 시작하여 10년 후에는 채택 속도에 따라 낮은 두 자릿수부터 시민 지정 시스템으로의 완전한 이행에 이르기까지. 실제 속도는 국가 서비스에 대한 자유 시장 대안이 나타나는 속도와 이행에 협력하려는 국가의 의지에 달려 있다. 시간의 함수는 선형일 필요가 없으며, 채택이 궁극적으로 도달할 지점에 상한을 설정할 이유는 없다.

% ============================================================
% 9. MIGRATION PATH
% ============================================================
\section{이행 경로}

이행은 세 단계를 거쳐 진행된다(그림~\ref{fig:migration}). \textbf{1단계: 오버레이}(RSN이 정보 계층으로 작동, 국가가 유일한 제공자), \textbf{2단계: 경쟁}(RSN이 기능적 대안 제공, 이중 시스템 사용), 그리고 \textbf{3단계: 이행}(RSN이 국가 등가물을 능가, 자발적 이주). 이행은 모든 단계에서 되돌릴 수 있다.

\begin{figure}[ht]
\centering
\begin{tikzpicture}[
    phase/.style={draw, rounded corners=3pt, minimum width=1.8cm, minimum height=0.8cm, align=center, font=\scriptsize, fill=black!5},
    arr/.style={-{Stealth[length=5pt]}, thick}
]
\node[phase] (p1) at (0,0) {\textbf{1단계}\\오버레이};
\node[phase] (p2) at (2.2,0) {\textbf{2단계}\\경쟁};
\node[phase] (p3) at (4.4,0) {\textbf{3단계}\\이행};

\draw[arr] (p1) -- (p2);
\draw[arr] (p2) -- (p3);
\draw[arr, dashed, gray] (p3.south) -- ++(0,-0.6) -| node[below,font=\tiny,pos=0.25] {복귀} (p2.south);
\end{tikzpicture}
\caption{3단계 이행---모든 단계에서 되돌릴 수 있음.}
\label{fig:migration}
\end{figure}

주된 압력 메커니즘은 재정적이다. 조건부 납세자가 ``경제적으로 유의미한 소수 $X$''---그에 대한 억압이 무시할 수 없는 경제적 피해를 초래하고 시민 불복종이 신빙성 있는 위협이 될 만큼 충분히 큰 집단---에 도달하면, 국가는 협상에 들어간다. 예시를 위해 $X \approx$ GDP의 $15\%$를 사용하지만, 실제 문턱은 특정 경제에 달려 있다.

% ============================================================
% 10. SECURITY ANALYSIS
% ============================================================
\section{보안 분석}

우리는 다섯 가지 적대자 범주를 고려한다. 개인 악의적 행위자, 조직된 집단, 기업 적대자, 국가 수준 적대자, 그리고 프로토콜 수준 적대자다.

\textbf{시빌(Sybil) 저항}~\cite{douceur2002}. 평판을 쌓는 것은 지속적이고 검증 가능한 상호작용을 요구한다. 성공적인 시빌 공격의 비용은 가짜 신원의 수에 따라 선형적으로, 목표 평판 수준에 따라 이차적으로 증가한다. 여러 DID를 병렬로 운영하는 것은 가능하지만 설계상 비싸다---각 신원은 진정한 활동을 통해 독립적으로 평판을 축적해야 한다. 자유 사회에서 이 비용은 안일한 남용을 억제한다. 독재 국가에서 병렬 DID는 구획화된 저항, 지하 조정, 더 안전한 암시장 항해를 가능하게 하는 생존 도구가 된다.

\textbf{공모 방어.} 폐쇄된 집단 내의 상호 보증 패턴은 소셜 그래프 분석을 통해 탐지 가능하다. 평판 집계 함수는 긴밀하게 군집된 출처의 주장을 할인한다.

\textbf{국가 수준 적대자.} 어니언 라우팅, 중앙집권적 인프라의 부재, 그리고 참여자 기반 전반에 걸친 검증자 역할의 분산은 억압이 어떤 이익과도 비례하지 않는 조치를 요구하도록 보장한다.

\textbf{프라이버시 대 책임성.} 가명성---익명성과 신원 공개 사이의 중간 지대---은 DID가 보유자의 현실 세계 신원을 드러내지 않고도 의미 있는 평판을 축적할 수 있게 한다.

% ============================================================
% 11. PRIVACY
% ============================================================
\section{프라이버시 고려사항}

RSN의 프라이버시 모델은 가명성에 기반한다. 보유자는 신원 공개를 통제한다. 최소(순수 가명), 선택적(특정 자격 증명 연결), 또는 완전(법적 신원 연결)이다. 게시된 주장은 영구적이다---시스템은 맥락적 정정은 지원하지만 삭제는 지원하지 않는다. 보유자는 손상된 DID를 버리고 축적한 평판을 포기하며 새로 시작할 수 있다.

% ============================================================
% 12. RELATED WORK
% ============================================================
\section{관련 연구}

W3C DID Core 명세~\cite{did2022}와 Ceramic Network~\cite{ceramic2023}는 신원 기반을 제공한다. EigenTrust~\cite{kamvar2003}는 중앙 권위 없는 분산 평판 집계를 입증했다. SBT~\cite{weyl2022}는 양도 불가능한 사회적 토큰을 도입했다. RSN은 품질 필터로서의 경제적 비용에 대한 강조와 급진주의에 가격을 매기는 메커니즘에서 차별화된다.

DAO~\cite{buterin2014}와 이차 투표~\cite{buterin2019}는 탈중앙화 거버넌스의 실현 가능성을 입증했다. RSN은 투표가 아니라 양방향 가격 협상으로부터 합의를 도출한다.

이 제안은 민간 서비스 제공에 관해 무정부 자본주의 이론~\cite{rothbard1973,hoppe2001}에 기대지만 두 가지 결정적 측면에서 벗어난다. 즉 합리성을 가정하기보다 앙심과 응보적 충동을 위해 설계하며, 혁명이 아니라 점진적 이행을 제안한다. 창발적 사회 계약이라는 개념은 하이에크의 자생적 질서(spontaneous order) 이론~\cite{hayek1973}에 선례를 둔다.

\textbf{무정부 아고리즘(anarcho-agorism).} RSN은 아고리스트 대항 경제학~\cite{konkin2006}과 상당한 공통 기반을 공유한다---특히 병렬 제도 구축과 자발적 교환에 대한 강조에서다. 그러나 아고리즘은 합리적 자기 이익을 충분한 조정 메커니즘으로 가정하는 경향이 있으며, 기존 국가 구조로부터의 구체적 이행 경로가 결여된 경우가 많다. RSN은 비합리적 행동 경향을 명시적으로 통합하고 점진적 이행을 위한 재정적 압력 메커니즘을 제공한다.

\textbf{실력주의(meritocracy).} RSN은 실력주의~\cite{young1958}가 구호가 아니라 시스템적 속성으로 어떻게 창발할 수 있는지에 대한 최초의 기능적 서술을 나타낼 수 있다. 전통적 실력주의 시스템은 ``실력''을 정의하고 집행할 중앙 권위를 요구하기 때문에 실패한다. RSN에서 실력은 양방향 평가로부터 창발한다. 가치를 제공하는 자는 평판을 축적하며, 어떤 위원회도 누가 실력을 가졌는지 결정하지 않는다.

\textbf{특권으로서의 재산.} RSN은 재산권을 자연적이고 절대적인 것으로 다루는 무정부 자본주의 및 오스트리아 학파의 취급에서 근본적으로 벗어난다. RSN 프레임워크에서 재산은 \emph{특권}---극단적인 경우 참여자가 공유된 규범을 근본적으로 위반할 때(배신, 실존적 위기에서 공동체를 방어하기를 거부, 체계적 착취) 공동체가 철회할 수 있는 사회적 인정---이다. 이것은 국가의 몰수가 아니라 양방향 관계 네트워크에 의한 사회적 인정의 철회다.

% ============================================================
% 13. LIMITATIONS
% ============================================================
\section{논의와 한계}

\textbf{콜드 스타트.} RSN의 가치는 네트워크 효과에 달려 있다. 초기 채택자는 참여가 문턱에 도달할 때까지 제한된 가치에 직면한다. 그러나 초기 단계부터도 DID 네트워크는 유용한 보충 정보원으로 기능한다. 초기 평판 평가와 권위 평가는 정교함이 증가하며 점진적으로 발전할 것으로 예상된다.

\textbf{반(反)무임승차와 접근 제어.} 대상 DID는 저평판 DID의 조회에 점진적 수수료와 접근 제한을 부과할 수 있다. 이는 이분법이 아니라 연속적 척도다. 조회하는 DID가 평판이 적을수록 받는 정보가 적고, 더 오래 기다리며, 더 많이 지불한다. 이 메커니즘은 수동적 소비보다 진정한 참여를 장려한다.

\textbf{접근의 불평등.} 주장을 게시하는 비용은 경제적으로 불리한 참여자의 정당한 불만을 걸러낼 수 있다. 완화책: 모든 참여자는 동시에 잠재적 검증자로 기능하며(또는 이 역할을 위임하며) 검증 수수료를 번다. 충분한 시간 지평에 걸쳐, 급진적이지 않은 참여자는 재정적 중립---검증 소득이 게시 비용을 대략 상쇄하는---에 접근해야 한다. 이는 통계적 기대이지 보장이 아니다.

\textbf{평판 불평등.} 초기 채택자는 후발 주자에게 장벽이 될 수 있는 이점을 축적한다. 그러나 평판 평가는 제3자 제공자가 수행하며, 이들은 자신의 집계 알고리즘을 통해 선점 이점을 완화하는 쪽을 선택할 수 있다. 좋은 평판으로 먼저 있는 것은 정당한 비교 경제적 이점이며---그것은 의도된 설계다.

\textbf{미해결 질문.} 최적 비용 함수, 집계 표준, 프로토콜 거버넌스, 그리고 AI 생성 합성 평판 데이터와의 상호작용은 향후 연구 영역으로 남는다.

본 논문은 RSN이 모든 상황에서 최적의 결과를 산출할 것이라고 주장하지 않는다. 본 논문은 RSN이 \emph{실현 가능한} 대안---기술적으로 구현 가능하고, 경제적으로 자기 지속적이며, 있는 그대로의 인간 행동 경향과 사회적으로 양립 가능한---을 나타낸다고 주장한다.

% ============================================================
% 14. CONCLUSION
% ============================================================
\section{결론}

본 논문은 가명적이고 검증 가능한 평판 교환을 가능하게 하는 탈중앙화 프로토콜인 평판 소셜 네트워크를 제시했다. 급진주의 비용화 원칙은 사기성 주장이 검열 없이 가격에서 배제되도록 보장한다. 제안된 이행 경로는 기존 기관으로부터의 점진적이고 되돌릴 수 있는 이행을 가능하게 한다. 이 설계는 이념적 변혁을 요구하기보다---옹졸함, 앙심, 응보적 만족에 대한 욕구를 포함하여---있는 그대로의 인간 본성을 통합한다. 그 결과는 유토피아가 아니라 정의에 접근할 수 있고, 평판이 획득되며, 부정행위의 비용이 그것을 초래한 당사자에 의해 부담되는 시스템이다.

% ============================================================
% REFERENCES
% ============================================================
\bibliography{references}

\end{document}
